% Source: https://artofproblemsolving.com/community/c3h1347637_graphing_caluculator_question
% Stable identity: aops-topic-1347637-q01
Write $p(x)=ax+b$. The two points give the system
\[
2a+b=5,\qquad 8a+b=1,
\qquad\text{with augmented matrix}\qquad
\begin{pmatrix}2&1&5\\8&1&1\end{pmatrix}.
\]
Row reduction: subtracting $4$ times row $1$ from row $2$ gives $(0,\,-3,\,-19)$, so $b=\dfrac{19}{3}$. Then $2a=5-\dfrac{19}{3}=-\dfrac{4}{3}$, so $a=-\dfrac{2}{3}$. The reduced row echelon form is
\[
\begin{pmatrix}1&0&-\dfrac{2}{3}\\[4pt]0&1&\dfrac{19}{3}\end{pmatrix},
\qquad\text{so}\qquad
\boxed{p(x)=-\dfrac{2}{3}x+\dfrac{19}{3}}.
\]
Check: $p(2)=\dfrac{-4+19}{3}=5$ and $p(8)=\dfrac{-16+19}{3}=1$.

\textbf{Why the degree condition.} Two points do not determine a polynomial of higher degree. For any number $c$, the polynomial $p(x)+c(x-2)(x-8)$ also passes through both points, so there are infinitely many such polynomials of degree $2$. In general, $n+1$ points with distinct $x$-coordinates determine a unique polynomial of degree at most $n$.
